% ============================================================
%  CHAPTER 10: DEPLOYMENT, LIMITATIONS, AND MAINTENANCE
% ============================================================

\chapter{DEPLOYMENT, LIMITATIONS, AND MAINTENANCE}
\label{chap:deployment}

\section{Production Deployment Architecture}
\label{sec:deployment_architecture}

Deploying QuantumSafe VPN involves establishing a hardened, resilient infrastructure consisting of two operational endpoints: the Linux VPN Gateway and the Android Mobile Endpoints. The gateway is deployed within an enterprise Demilitarized Zone (DMZ) or hosted on a virtual private server (VPS) with public WAN connectivity.

\section{Linux Gateway Server Provisioning and Network Configuration}
\label{sec:linux_server_setup}

The gateway operates on Kali Linux 2024.x, Ubuntu Server 22.04 LTS, or Debian 12 (Kernel 5.15+).

\subsection{Kernel IPv4 Forwarding and NAT Masquerading}
The Linux kernel routes packets between the virtual \texttt{tun0} interface and physical \texttt{eth0} WAN interface via sysctl and \texttt{iptables}:
\begin{lstlisting}[language=bash, basicstyle=\ttfamily\scriptsize]
# Enable runtime and persistent kernel IPv4 forwarding
sudo sysctl -w net.ipv4.ip_forward=1
echo "net.ipv4.ip_forward = 1" | sudo tee -a /etc/sysctl.d/99-vpn-forward.conf

# Masquerade outbound traffic from VPN clients (10.8.0.0/24) via physical WAN (eth0)
sudo iptables -t nat -A POSTROUTING -s 10.8.0.0/24 -o eth0 -j MASQUERADE
sudo iptables -A FORWARD -i tun0 -o eth0 -j ACCEPT
sudo iptables -A FORWARD -i eth0 -o tun0 -m state --state RELATED,ESTABLISHED -j ACCEPT
sudo iptables -A INPUT -p tcp --dport 7000 -j ACCEPT
\end{lstlisting}

\subsection{Automated Daemon Startup Scripting}
The gateway daemon is launched via \texttt{scripts/start\_server.sh}:
\begin{lstlisting}[language=bash, basicstyle=\ttfamily\scriptsize]
#!/usr/bin/env bash
set -e
if [ "$EUID" -ne 0 ]; then echo "[-] Run as root"; exit 1; fi
bash "$(dirname "$0")/setup_network.sh"
export PYTHONPATH="." VPN_SERVER_PORT=7000
python3 vpn/tunnel_server.py
\end{lstlisting}

\section{Android Client Deployment and Sideloading}
\label{sec:android_client_deployment}

The mobile client is installed on Android 8.0--14 devices via ADB sideloading or direct package installation:
\begin{lstlisting}[language=bash, basicstyle=\ttfamily\scriptsize]
adb install -r dist/QuantumSafe-VPN-release.apk
\end{lstlisting}
Upon first launch, the user configures the target gateway IP (\texttt{11.12.5.8}) and port (\texttt{7000}) in the Server Config tab, returns to Home, and taps Connect, accepting the system-level \texttt{VpnService} permission prompt.

\section{Operational Troubleshooting Runbook}
\label{sec:troubleshooting_runbook}

Table~\ref{tab:troubleshooting_guide} details common operational issues, diagnostic indicators, and verified resolution procedures.

\begin{table}[htbp]
  \centering
  \scriptsize
  \setlength{\tabcolsep}{3pt}
  \caption{Operational Troubleshooting and Diagnostic Runbook.}
  \label{tab:troubleshooting_guide}
  \begin{tabularx}{\linewidth}{@{}p{3.2cm}p{3.8cm}Y@{}}
    \toprule
    \textbf{Observed Symptom} & \textbf{Probable Root Cause} & \textbf{Prescribed Resolution Procedure} \\
    \midrule
    Connection Timeout (UI ERROR) & Port 7000 blocked by firewall; daemon inactive & 1. Verify daemon: \texttt{ss -tulpn | grep 7000}.\\ & & 2. Allow TCP 7000 in firewall.\\ & & 3. Confirm network route to gateway IP. \\
    Connected; No Internet Access & Kernel IP forwarding off; NAT rule missing & 1. Check: \texttt{sysctl net.ipv4.ip\_forward}.\\ & & 2. Verify NAT: \texttt{iptables -t nat -L -v -n}.\\ & & 3. Run \texttt{sudo bash scripts/setup\_network.sh}. \\
    Signature Failure & Server key altered; MitM active & 1. Inspect server log for key renewal.\\ & & 2. Clear cached key in app Settings.\\ & & 3. Re-verify \texttt{server\_mldsa65\_public.key}. \\
    UI Freezes upon Connect & Socket loop bypass failed & 1. Confirm \texttt{protect(socket)} returned true.\\ & & 2. Ensure OS battery manager allows app.\\ & & 3. Restart application process. \\
    \bottomrule
  \end{tabularx}
\end{table}

\section{System Limitations and Maintenance Lifecycle}
\label{sec:system_limitations}

\begin{enumerate}
    \item TCP Encapsulation Overhead: Tunneling IP packets over TCP can induce throughput degradation (TCP-over-TCP meltdown) on cellular links experiencing $> 5\%$ packet loss. A UDP/DTLS transport engine with automated TCP fallback is planned for subsequent iterations.
    \item Single-Hop Star Topology: The system connects clients to a centralized gateway; decentralized peer-to-peer mesh routing is not currently supported.
    \item IPv4 Exclusivity: The virtual TUN interface captures IPv4 traffic; native IPv6 dual-stack encapsulation remains an ongoing development objective.
    \item Pure-Kotlin Mathematical Primitives: While cross-platform and memory-safe, pure-Kotlin polynomial arithmetic does not leverage platform SIMD vector registers (ARM NEON), resulting in slightly higher CPU cycle counts than optimized C/assembly.
\end{enumerate}

The server enforces zero-logging (discarding packet payloads immediately and keeping session tables in volatile RAM), and server ML-DSA identity keys are rotated annually.
